\chapter*{参考文献}
\addcontentsline{toc}{chapter}{参考文献}
\markboth{参考文献}{参考文献}
\begingroup
\small\raggedright
\begin{list}{}{\leftmargin=2.1em\labelwidth=0pt\labelsep=0pt
 \itemindent=-2.1em\parsep=0pt\itemsep=3pt}
\item S. Acid and L. M. de Campos. Searching for Bayesian network structures in the space of restricted acyclic partially directed graphs. Journal of Artificial Intelligence Research, 18:445--490, 2003.
\item J. Aldrich. Autonomy. Oxford Economic Papers, 41:15--34, 1989.
\item R. A. Ali, T. S. Richardson, and P. Spirtes. Markov equivalence for ancestral graphs. The Annals of Statistics, 37:2808--2837, 2009.
\item E. S. Allman, C. Matias, and J. A. Rhodes. Identifiability of parameters in latent structure models with many observed variables. The Annals of Statistics, 37: 3099--3132, 2009.
\item N. Ancona, D. Marinazzo, and S. Stramaglia. Radial basis function approach to nonlinear Granger causality of time series. Physical Review E, 70:056221--1--7, 2004.
\item S. A. Andersson, D. Madigan, and M. D. Perlman. Alternative Markov property for chain graphs. Scandinavian Journal of Statistics, 28:33--86, 2001.
\item A. Aspect, P. Grangier, and G. Roger. Experimental tests of realistic local theories via Bell's theorem. Physical Review Letters, 47:460--467, 1981.
\item N. Ay and D. Polani. Information flows in causal networks. Advances in Complex Systems, 11(1):17--41, 2008.
\item R. R. Baldwin, W. E. Cantey, H. Maisel, and J. P. McDermott. The optimum strategy in blackjack. Journal of the American Statistical Association, 51(275): 429--439, 1956.
\item A. Balke. Probabilistic Counterfactuals: Semantics, Computation, Applications. PhD thesis, University of California, Los Angeles, CA, 1995.
\item A. Balke and J. Pearl. Bounds on treatment effects from studies with imperfect compliance. Journal of the American Statistical Association, 92:1172--1176, 1997.
\item E. Bareinboim and J. Pearl. Transportability from multiple environments with limited experiments: Completeness results. In Advances in Neural Information Processing Systems 27 (NIPS), pages 280--288, 2014.
\item E. Bareinboim and J. Pearl. Causal inference from big data: Theoretical foundations and the data-fusion problem. Proceedings of the National Academy of Sciences, 113(27):7345--7352, 2016.
\item E. Bareinboim, A. Forney, and J. Pearl. Bandits with unobserved confounders: A causal approach. In Advances in Neural Information Processing Systems 28 (NIPS), pages 1342--1350, 2015.
\item S. Bauer, B. Schölkopf, and J. Peters. The arrow of time in multivariate time series. In Proceedings of the 33rd International Conference on Machine Learning (ICML), pages 2043--2051, 2016.
\item J. Bell. On the Einstein-Podolsky-Rosen paradox. Physics, 1:195--200, 1964. Reprinted in J. Bell: Speakable and unspeakable in quantum mechanics, Cambridge University Press, Cambridge, UK, 1987.
\item R. Bellman. A Markovian decision process. Indiana University Mathematics Journal, 6:679--684, 1957.
\item S. Ben-David, T. Lu, T. Luu, and D. Pál. Impossibility theorems for domain adaptation. In Proceedings of the International Conference on Artificial Intelligence and Statistics 13 (AISTATS), pages 129--136, 2010.
\item C. Bennett. The thermodynamics of computation --- a review. International Journal of Theoretical Physics, 21:905--940, 1982.
\item J. O. Berger. Statistical Decision Theory and Bayesian Analysis. Springer, New York, NY, 1985.
\item J. Berkson. Limitations of the application of fourfold table analysis to hospital data. Biometrics Bulletin, 2:47--53, 1946.
\item M. Besserve, N. Shajarisales, D. Janzing, and B. Schölkopf. Causal inference through spectral independence in linear dynamical systems. in preparation.
\item A. Bhatt. Evolution of clinical research: A history before and beyond James Lind. Perspectives in Clinical Research, 1(1):6--10, 2010.
\item K. A. Bollen. Structural Equations with Latent Variables. Wiley, New York, NY, 1989.
\item B. Bonet. Instrumentality tests revisited. In Proceedings of the 17th Conference on Uncertainty in Artificial Intelligence (UAI), pages 48--55, 2001.
\item S. Bongers, J. Peters, B. Schölkopf, and J. M. Mooij. Structural causal models: Cycles, marginalizations, exogenous reparametrizations and reductions. ArXiv e-prints (1611.06221), 2016.
\item L. Bottou, J. Peters, J. Quiñonero-Candela, D. X. Charles, D. M. Chickering, E. Portugualy, D. Ray, P. Simard, and E. Snelson. Counterfactual reasoning and learning systems: The example of computational advertising. Journal of Machine Learning Research, 14:3207--3260, 2013.
\item R. J. Bowden and D. A. Turkington. Instrumental Variables. Econometric Society Monographs. Cambridge University Press, New York, NY, 1990.
\item C. Brito and J. Pearl. A new identification condition for recursive models with correlated errors. Structural Equation Modeling, 9:459--474, 2002a.
\item C. Brito and J. Pearl. Generalized instrumental variables. In Proceedings of the 18th Conference on Uncertainty in Artificial Intelligence (UAI), pages 85--93, 2002b.
\item P. J. Brockwell and R. A. Davis. Time Series: Theory and Methods. Springer, New York, NY, 2nd edition, 1991.
\item P. Bühlmann and S. A. van de Geer. Statistics for High-Dimensional Data: Methods, Theory and Applications. Springer Series in Statistics. Springer, New York, NY, 2011.
\item P. Bühlmann, J. Peters, and J. Ernest. CAM: Causal additive models, high-dimensional order search and penalized regression. The Annals of Statistics, 42(6):2526--2556, 2014.
\item W. L. Buntine. Theory refinement on Bayesian networks. In In Proceedings of the 7th Annual Conference on Uncertainty in Artificial Intelligence (UAI), pages 52--60, 1991.
\item R. M. J. Byrne. The Rational Imagination: How People Create Alternatives to Reality. MIT Press, Cambridge, MA, 2007.
\item G. Chaitin. On the length of programs for computing finite binary sequences. Journal of the Assocation for Computing Machinery, 13(4):547--569, 1966.
\item O. Chapelle, B. Schölkopf, and A. Zien. Semi-Supervised Learning. MIT Press, Cambridge, MA, 2006.
\item C. R. Charig, D. R. Webb, S. R. Payne, and J. E. A. Wickham. Comparison of treatment of renal calculi by open surgery, percutaneous nephrolithotomy, and extracorporeal shockwave lithotripsy. British Medical Journal (Clin Res Ed), 292(6254):879--882, 1986.
\item R. Chaves, L. Luft, T. O. Maciel, D. Gross, D. Janzing, and B. Schölkopf. Inferring latent structures via information inequalities. In Proceedings of the 30th Annual Conference on Uncertainty in Artificial Intelligence (UAI), pages 112--121, 2014.
\item D. M. Chickering. Learning Bayesian networks is NP-complete. In Learning from Data: Artificial Intelligence and Statistics V, pages 121--130. Springer, New York, NY, 1996.
\item D. M. Chickering. Optimal structure identification with greedy search. Journal of Machine Learning Research, 3:507--554, 2002.
\item T. Claassen, J. M. Mooij, and T. Heskes. Learning sparse causal models is not NP-hard. In Proceedings of the 29th Annual Conference on Uncertainty in Artificial Intelligence (UAI), pages 172--181, 2013.
\item J. Clauser, M. Horne, A. Shimony, and R. Holt. Proposed experiment to test local hidden-variable theories. Physical Review Letters, 23:880--884, 1969.
\item D. Colombo, M. H. Maathuis, M. Kalisch, and T. S. Richardson. Learning high-dimensional directed acyclic graphs with latent and selection variables. The Annals of Statistics, 40:294--321, 2012.
\item J. Comley and D. Dowe. General Bayesian networks and asymmetric languages. In Proceedings of the 2nd Hawaii International Conference on Statistics and Related Fields, 2003.
\item P. Comon. Independent component analysis --- a new concept? Signal Processing, 36:287--314, 1994.
\item S. A. Cook. The complexity of theorem-proving procedures. In Proceedings of the 3rd Annual ACM Symposium on Theory of Computing, pages 151--158, 1971.
\item G. Cooper and C. Yoo. Causal discovery from a mixture of experimental and observational data. In Proceedings of the 15th Annual Conference on Uncertainty in Artificial Intelligence (UAI), pages 116--125, 1999.
\item D. Corfield, B. Schölkopf, and V. Vapnik. Falsificationism and statistical learning theory: Comparing the Popper and Vapnik-Chervonenkis dimensions. Journal for General Philosophy of Science, 40(1):51--58, 2009.
\item T. Cover and J. Thomas. Elements of Information Theory. Wiley Series in Telecommunications and Signal Processing, Wiley, New York, NY, 1991.
\item D. R. Cox. Planning of Experiments. Wiley, New York, NY, 1958.
\item J. Cussens. Bayesian network learning with cutting planes. In Proceedings of the 27th Annual Conference on Uncertainty in Artificial Intelligence (UAI), pages 153--160, 2011.
\item P. Daniušis, D. Janzing, J. M. Mooij, J. Zscheischler, B. Steudel, K. Zhang, and B. Schölkopf. Inferring deterministic causal relations. In Proceedings of the 26th Annual Conference on Uncertainty in Artificial Intelligence (UAI), pages 143--150, 2010.
\item D. Danks and S. Plis. Learning causal structure from undersampled time series, 2013. URL \url{http://repository.cmu.edu/cgi/viewcontent.cgi?article=1638&context=philosophy}. Results were presented at NIPS 2013 workshop on causality; last visit of website: 31.01.2017.
\item G. Darmois. Analyse générale des liaisons stochastiques. Revue de l'Institut International de Statistique, 21:2--8, 1953.
\item D. Dash. Restructing dynamic causal systems in equilibrium. In Proceedings of the 10th International Conference on Artificial Intelligence and Statistics (AIS-TATS), pages 81--88, 2005.
\item A. P. Dawid. Conditional independence in statistical theory. Journal of the Royal Statistical Society, Series B: Statistical Methodology (with discussion), 41(1): 1--31, 1979.
\item A. P. Dawid. Statistical causality from a decision-theoretic perspective. Annual Review of Statistics and Its Application, 2:273--303, 2015.
\item C. P. De Campos and Q. Ji. Efficient structure learning of Bayesian networks using constraints. Journal of Machine Learning Research, 12:663--689, 2011.
\item L. Devroye, L. Györfi, and G. Lugosi. A Probabilistic Theory of Pattern Recognition, volume 31 of Applications of Mathematics. Springer, New York, NY, 1996.
\item V. Didelez, S. Meng, and N. A. Sheehan. Assumptions of IV methods for observational epidemiology. Statistical Science, 25:22--40, 2010.
\item C. Diks and V. Panchenko. A new statistic and practical guidelines for nonparametric Granger causality testing. Journal of Economic Dynamics and Control, 30(9--10):1647--1669, 2006.
\item D. DiVincenzo. Two-qubit gates are universal for quantum computation. Physical Review A, 51:1015--1022, 1995.
\item M. Drton, M. Eichler, and T. S. Richardson. Computing maximum likelihood estimates in recursive linear models with correlated errors. Journal of Machine Learning Research, 10:2329--2348, 2009a.
\item M. Drton, B. Sturmfels, and S. Sullivant. Lectures on Algebraic Statistics, volume 39 of Oberwolfach Seminars. Birkhäuser, Basel, 2009b.
\item M. Druzdzel and H. Simon. Causality in Bayesian belief networks. In In Proceedings of the 9th Annual Conference on Uncertainty in Artificial Intelligence (UAI), pages 3--11, 1993.
\item M. J. Druzdzel and H. van Leijen. Causal reversibility in Bayesian networks. Journal of Experimental and Theoretical Artificial Intelligence, 13(1):45--62, 2001.
\item R. M. Dudley. Real Analysis and Probability. Cambridge University Press, New York, NY, 2002.
\item D. Eaton and K. P. Murphy. Exact Bayesian structure learning from uncertain interventions. In Proceedings of the 11th International Conference on Artificial Intelligence and Statistics (AISTATS), pages 107--114, 2007.
\item F. Eberhardt and R. Scheines. Interventions and causal inference. Philosophy of Science, 74(5):981--995, 2007.
\item F. Eberhardt, C. Glymour, and R. Scheines. On the number of experiments sufficient and in the worst case necessary to identify all causal relations among n variables. In Proceedings of the 21st Annual Conference on Uncertainty in Artificial Intelligence (UAI), pages 178--184, 2005.
\item F. Eberhardt, P. O. Hoyer, and R. Scheines. Combining experiments to discover linear cyclic models with latent variables. In Proceedings of the 13th International Conference on Artificial Intelligence and Statistics (AISTATS), pages 185--192, 2010.
\item M. Eichler. Graphical modelling of multivariate time series. Probability Theory and Related Fields, pages 1--36, 2011.
\item M. Eichler. Causal inference in time series analysis. In C. Berzuini, P. Dawid, and L. Bernardinelli, editors, Causality: Statistical Perspectives and Applications, pages 327--354. Wiley, Chichester, UK, 2012.
\item J. Ellenberg. How Not to Be Wrong: The Power of Mathematical Thinking. Penguin Press, London, UK, 2014.
\item R. F. Engle, D. F. Hendry, and J.-F. Richard. Exogeneity. Econometrica, 51(2): 277--304, 1983.
\item R. J. Evans. Graphical methods for inequality constraints in marginalized DAGs. In Proceedings of the 22nd Workshop on Machine Learning and Signal Processing, pages 1--12, 2012.
\item R. J. Evans. Margins of discrete Bayesian networks. ArXiv e-prints (arXiv:1501.02103), 2015.
\item R. J. Evans and T. S. Richardson. Markovian acyclic directed mixed graphs for discrete data. The Annals of Statistics, 42(2):1452--1482, 2014.
\item T. S. Ferguson. Mathematical Statistics --- A Decision Theoretic Approach. Academic Press, New York, NY and London, UK, 1967.
\item A. Fink. The binomial ideal of the intersection axiom for conditional probabilities. Journal of Algebraic Combinatorics, 33(3):455--463, 2011.
\item R. A. Fisher. Statistical Methods for Research Workers. Oliver \& Boyd, Edinburgh, UK, 1925.
\item W. T. Freeman. The generic viewpoint assumption in a framework for visual perception. Nature, 368(6471):542--545, 1994.
\item R. Frisch and F. V. Waugh. Partial time regressions as compared with individual trends. Econometrica, 1(4):387--401, 1933.
\item R. Frisch, T. Haavelmo, T. C. Koopmans, and J. Tinbergen. Autonomy of Economic Relations. Series: Memorandum fra Universitets Socialøkonomiske Institutt. Universitets Socialøkonomiske Institutt, Oslo, Norway, 1948.
\item K. Friston, L. Harrison, and W. Penny. Dynamic causal modelling. NeuroImage, 19:1273--1302, 2003.
\item K. Friston, J. Daunizeau, and K. Stephan. Model selection and gobbledygook: Response to Lohmann et al. NeuroImage, 75:275--278, 2013.
\item M. Frydenberg. The chain graph Markov property. Scandinavian Journal of Statistics, 17(4):333--353, 1990.
\item K. Fukumizu, A. Gretton, X. Sun, and B. Schölkopf. Kernel measures of conditional dependence. In Advances in Neural Information Processing Systems 20 (NIPS), pages 489--496, 2008.
\item J. A. Gagnon-Bartsch and T. P. Speed. Using control genes to correct for unwanted variation in microarray data. Biostatistics, 13:539--552, 2012.
\item D. Galles and J. Pearl. An axiomatic characterization of causal counterfactuals. Foundations of Science, 3(1):151--182, 1998.
\item D. Geiger and D. Heckerman. Learning Gaussian networks. In Proceedings of the 10th Annual Conference on Uncertainty in Artificial Intelligence (UAI), pages 235--243, 1994a.
\item D. Geiger and D. Heckerman. Learning Bayesian networks: The combination of knowledge and statistical data. In Proceedings of the 10th Annual Conference on Uncertainty in Artificial Intelligence (UAI), pages 293--301, 1994b.
\item D. Geiger and C. Meek. Graphical models and exponential families. In Proceedings of 14th Annual Conference on Uncertainty in Artificial Intelligence (UAI), pages 156--165, 1998.
\item P. Geiger, D. Janzing, and B. Schölkopf. Estimating causal effects by bounding confounding. In Proceedings of the 30th Conference on Uncertainty in Artificial Intelligence (UAI), pages 240--249, 2014.
\item D. Geradin and I. Girgenson. The counterfactual method in EU competition law: The cornerstone of the effects-based approach. Available at SSRN: \url{http://ssrn.com/abstract=1970917}, 2011.
\item C. W. J. Granger. Testing for causality: A personal viewpoint. Journal of Economic Dynamics and Control, 2(1):329--352, 1980.
\item C. W. J. Granger. Some recent development in a concept of causality. Journal of Econometrics, 39(1--2):199--211, 1988.
\item A. Gretton, K. Fukumizu, C. H. Teo, L. Song, B. Schölkopf, and A. Smola. A kernel statistical test of independence. In Advances in Neural Information Processing Systems 20 (NIPS), pages 585--592, 2008.
\item P. D. Grünwald. The Minimum Description Length Principle. MIT Press, Cambridge, MA, 2007.
\item I. Guyon. Challenge: Cause-effect pairs, 2013. URL \url{https://www.kaggle.com/c/cause-effect-pairs/}. Results were presented at NIPS 2013 workshop \url{http://clopinet.com/isabelle/Projects/NIPS2013/}; last visit of websites: 19.07.2016.
\item J. Gwiazda, E. Ong, R. Held, and F. Thorn. Vision: Myopia and ambient night-time lighting. Nature, 404:144, 2000.
\item T. Haavelmo. The statistical implications of a system of simultaneous equations. Econometrica, 11(1):1--12, 1943.
\item T. Haavelmo. The probability approach in econometrics. Econometrica, 12:S1--S115 (supplement), 1944.
\item J. Y. Halpern. Axiomatizing causal reasoning. Journal of Artificial Intelligence Research, 12:317--337, 2000.
\item J. Y. Halpern. Actual Causality. MIT Press, Cambridge, MA, 2016.
\item N. R. Hansen and A. Sokol. Causal interpretation of stochastic differential equations. Electronic Journal of Probability, 19(100):1--24, 2014.
\item T. Hastie, R. Tibshirani, and J. Friedman. The Elements of Statistical Learning: Data Mining, Inference and Prediction. Springer, New York, NY, 2nd edition, 2009.
\item D. M. A. Haughton. On the choice of a model to fit data from an exponential family. The Annals of Statistics, 16(1):342--355, 1988.
\item A. Hauser and P. Bühlmann. Characterization and greedy learning of interventional Markov equivalence classes of directed acyclic graphs. Journal of Machine Learning Research, 13:2409--2464, 2012.
\item A. Hauser and P. Bühlmann. Two optimal strategies for active learning of causal models from interventional data. International Journal of Approximate Reasoning, 55:926--939, 2014.
\item A. Hauser and P. Bühlmann. Jointly interventional and observational data: Estimation of interventional Markov equivalence classes of directed acyclic graphs. Journal of the Royal Statistical Society, Series B: Statistical Methodology, 77: 291--318, 2015.
\item D. M. Hausman and J. Woodward. Independence, invariance and the causal Markov condition. The British Society for the Philosophy of Science, 50:521--583, 1999.
\item D. Heckerman, C. Meek, and G. Cooper. A Bayesian approach to causal discovery. In C. Glymour and G. Cooper, editors, Computation, Causation, and Discovery, pages 141--165. MIT Press, Cambridge, MA, 1999.
\item R. Hemmecke, S. Linder, and M. Studený. Characteristic imsets for learning Bayesian network structure. International Journal of Approximate Reasoning, 53:1336--1349, 2012.
\item M. A. Hernán and J. M. Robins. Instruments for causal inference: An epidemiologists dream? Epidemiology, 17:360--372, 2006.
\item P. W. Holland. Statistics and causal inference. Journal of the American Statistical Association, 81:968--970, 1986.
\item K. D. Hoover. Causality in economics and econometrics. In S. N. Durlauf and L. E. Blume, editors, The New Palgrave Dictionary of Economics. Palgrave Macmillan, Basingstoke, UK, 2nd edition, 2008.
\item D. G. Horvitz and D. J. Thompson. A generalization of sampling without replacement from a finite universe. Journal of the American Statistical Association, 47(260):663--685, 1952.
\item P. O. Hoyer, A. Hyvärinen, R. Scheines, P. Spirtes, J. Ramsey, G. Lacerda, and S. Shimizu. Causal discovery of linear acyclic models with arbitrary distributions. In Proceedings of the 24th Annual Conference on Uncertainty in Artificial Intelligence (UAI), pages 282--289, 2008a.
\item P. O. Hoyer, S. Shimizu, A. J. Kerminen, and M. Palviainen. Estimation of causal effects using linear non-Gaussian causal models with hidden variables. International Journal of Approximate Reasoning, 49(2):362--378, 2008b.
\item P. O. Hoyer, D. Janzing, J. M. Mooij, J. Peters, and B. Schölkopf. Nonlinear causal discovery with additive noise models. In Advances in Neural Information Processing Systems 21 (NIPS), pages 689--696, 2009.
\item Y. Huang and M. Valtorta. Pearl's calculus of intervention is complete. In Proceedings of the 22nd Annual Conference on Uncertainty in Artificial Intelligence (UAI), pages 217--224, 2006.
\item L Hurwicz. On the structural form of interdependent systems. In E. Nagel, P. Suppes, and A. Tarski, editors, Logic, Methodology and Philosophy of Science, Proceedings of the 1960 International Congress, pages 232--239. Stanford University Press, Stanford, CA, 1962.
\item A. Hyttinen, F. Eberhardt, and P. O. Hoyer. Learning linear cyclic causal models with latent variables. Journal of Machine Learning Research, 13(1):3387--3439, 2012.
\item A. Hyttinen, P. O. Hoyer, F. Eberhardt, and M. Järvisalo. Discovering cyclic causal models with latent variables: A general SAT-based procedure. In Proceedings of the 29th Annual Conference on Uncertainty in Artificial Intelligence (UAI), pages 301--310, 2013.
\item A. Hyttinen, S. Plis, M. Järvisalo, F. Eberhardt, and D. Danks. Causal discovery from subsampled time series data by constraint optimization. In Proceedings of the 8th International Conference on Probabilistic Graphical Models (PGM), pages 216--227, 2016.
\item A. Hyvärinen, J. Karhunen, and E. Oja. Independent Component Analysis. Adaptive and Learning Systems for Signal Processing, Communications, and Control. Wiley, New York, NY, 2001.
\item A. Hyvärinen, S. Shimizu, and P. Hoyer. Causal modelling combining instantaneous and lagged effects: An identifiable model based on non-Gaussianity. In Proceedings of the 25th International Conference on Machine Learning (ICML), pages 424--431, 2008.
\item G. W. Imbens and J. Angrist. Identification and estimation of local average treatment effects. Econometrica, 62(2):467--75, 1994.
\item G. W. Imbens and D. B. Rubin. Discussion of: C̀̀ausal Diagrams for Empirical Research'' by J. Pearl. Biometrika, 82(4):694--695, 1995.
\item G. W. Imbens and D. B. Rubin. Causal Inference for Statistics, Social, and Biomedical Sciences: An Introduction. Cambridge University Press, New York, NY, 2015.
\item T. Jaakkola, D. Sontag, A. Globerson, and M. Meila. Learning Bayesian network structure using LP relaxations. Proceedings of the 13th International Conference on Artificial Intelligence and Statistics (AISTATS), pages 358--365, 2010.
\item L. Jacob, J. A. Gagnon-Bartsch, and T. P. Speed. Correcting gene expression data when neither the unwanted variation nor the factor of interest are observed. Biostatistics, 17(1):16--28, 2016.
\item D. Janzing. On the entropy production of time series with unidirectional linearity. Journal of Statistical Physics, 138:767--779, 2010.
\item D. Janzing and B. Schölkopf. Causal inference using the algorithmic Markov condition. IEEE Transactions on Information Theory, 56(10):5168--5194, 2010.
\item D. Janzing and B. Schölkopf. Semi-supervised interpolation in an anticausal learning scenario. Journal of Machine Learning Research, 16:1923--1948, 2015.
\item D. Janzing and B. Steudel. Justifying additive-noise-based causal discovery via algorithmic information theory. Open Systems and Information Dynamics, 17(2):189--212, 2010.
\item D. Janzing, J. Peters, J. M. Mooij, and B. Schölkopf. Identifying confounders using additive noise models. In Proceedings of the 25th Annual Conference on Uncertainty in Artificial Intelligence (UAI), pages 249--257, 2009a.
\item D. Janzing, X. Sun, and B. Schölkopf. Distinguishing cause and effect via second order exponential models. ArXiv e-prints (0910.5561), 2009b.
\item D. Janzing, P. O. Hoyer, and B. Schölkopf. Telling cause from effect based on high-dimensional observations. In Proceedings of the 27th International Conference on Machine Learning (ICML), pages 479--486, 2010.
\item D. Janzing, E. Sgouritsa, O. Stegle, J. Peters, and B. Schölkopf. Detecting low-complexity unobserved causes. In Proceedings of the 27th Annual Conference on Uncertainty in Artificial Intelligence (UAI), pages 383--391, 2011.
\item D. Janzing, J. M. Mooij, K. Zhang, J. Lemeire, J. Zscheischler, P. Daniušis, B. Steudel, and B. Schölkopf. Information-geometric approach to inferring causal directions. Artificial Intelligence, 182--183:1--31, 2012.
\item D. Janzing, D. Balduzzi, M. Grosse-Wentrup, and B. Schölkopf. Quantifying causal influences. The Annals of Statistics, 41(5):2324--2358, 2013.
\item D. Janzing, B. Steudel, N. Shajarisales, and B. Schölkopf. Justifying information-geometric causal inference. In V. Vovk, H. Papadopolous, and A. Gammerman, editors, Measures of Complexity, Festschrift for Alexey Chervonencis, pages 253--265. Springer, Heidelberg, Germany, 2015.
\item D. Janzing, R. Chaves, and B. Schölkopf. Algorithmic independence of initial condition and dynamical law in thermodynamics and causal inference. New Journal of Physics, 18(093052):1--13, 2016.
\item M. Kääriäinen. Generalization error bounds using unlabeled data. In Proceedings of the 18th Annual Conference on Learning Theory (COLT), pages 127--142, 2005.
\item M. Kalisch and P. Bühlmann. Estimating high-dimensional directed acyclic graphs with the PC-algorithm. Journal of Machine Learning Research, 8:613--636, 2007.
\item M. Kalisch, M. Mächler, D. Colombo, M. H. Maathuis, and P. Bühlmann. Causal inference using graphical models with the R package pcalg. Journal of Statistical Software, 47(11):1--26, 2012.
\item C. Kang and J. Tian. Inequality constraints in causal models with hidden variables. In Proceedings of the 22th Annual Conference on Uncertainty in Artificial Intelligence (UAI), pages 233--240, 2006.
\item Y. Kano and S. Shimizu. Causal inference using nonnormality. In Proceedings of the International Symposium on the Science of Modeling, the 30th Anniversary of the Information Criterion, pages 261--270, 2003.
\item A. Kela, K. von Prillwitz, J. ˚Aberg, R. Chaves, and D. Gross. Semidefinite tests for latent causal structures. ArXiv e-prints (1701.00652), 2017.
\item M. Koivisto. Advances in exact Bayesian structure discovery in Bayesian networks. In Proceedings of the 22nd Annual Conference on Uncertainty in Artificial Intelligence (UAI), pages 241--248, 2006.
\item M. Koivisto and K. Sood. Exact Bayesian structure discovery in Bayesian networks. Journal of Machine Learning Research, 5:549--573, 2004.
\item D. Koller and N. Friedman. Probabilistic Graphical Models: Principles and Techniques. MIT Press, Cambridge, MA, 2009.
\item A. Kolmogorov. Three approaches to the quantitative definition of information. Problems of Information Transmission, 1(1):3--11, 1965.
\item T. C. Koopmans. When is an equation system complete for statistical purposes? In T. C. Koopmans, editor, Statistical Inference in Dynamic Economic Models, pages 393--409. Wiley and Chapman \& Hall, New York, NY, and London, UK, 1950.
\item K. Korb, L. Hope, A. Nicholson, and K. Axnick. Varieties of causal intervention. In Proceedings of the Pacific Rim Conference on Artifical Intelligence, pages 322--331, 2004.
\item S. Kpotufe, E. Sgouritsa, D. Janzing, and B. Schölkopf. Consistency of causal inference under the additive noise model. In Proceedings of the 31st International Conference on Machine Learning, ICML 2014, Beijing, China, pages 478--486, 2014.
\item G. Lacerda, P. Spirtes, J. Ramsey, and P. O. Hoyer. Discovering cyclic causal models by independent components analysis. In Proceedings of the 24th Annual Conference on Uncertainty in Artificial Intelligence (UAI), pages 366--374, 2008.
\item S. L. Lauritzen. Graphical Models. Oxford University Press, New York, NY, 1996.
\item S. L. Lauritzen. Discussion on causality. Scandinavian Journal of Statistics, 31(2): 189--193, 2004.
\item S. L. Lauritzen and T. S. Richardson. Chain graph models and their causal interpretations. Journal of the Royal Statistical Society, Series B: Statistical Methodology (with discussion), 64(3):321--361, 2002.
\item S. L. Lauritzen and N. Wermuth. Graphical models for associations between variables, some of which are qualitative and some quantitative. Annals of Statistics, 17(1):31--57, 1989.
\item J. Lemeire and E. Dirkx. Causal models as minimal descriptions of multivariate systems. \url{http://parallel.vub.ac.be/~jan/}, 2006.
\item J. Lemeire and D. Janzing. Replacing causal faithfulness with algorithmic independence of conditionals. Minds and Machines, 23:227--249, 2013.
\item L. A. Levin. Universal sequential search problems. Problems of Information Transmission, 9(3):115--116, 1973. (Translated into English by B. A. Trakhtenbrot: `À survey of Russian approaches to perebor (brute-force searches) algorithms,'' Annals of the History of Computing 6(4): 384--400, 1984).
\item M. Levine, D. R. Hunter, and D. Chauveau. Maximum smoothed likelihood for multivariate mixtures. Biometrika, 98(2):403--416, 2011.
\item M. Li and P. Vitányi. An Introduction to Kolmogorov Complexity and Its Applications. Springer, New York, NY, 3rd edition, 1997.
\item G. Lohmann, K. Erfurth, K. Müller, and R. Turner. Critical comments on dynamic causal modelling. NeuroImage, 59:2322--2329, 2012a.
\item G. Lohmann, K. Müller, and R. Turner. Response to commentaries on our paper: Critical comments on dynamic modeling. NeuroImage, 75:279--281, 2012b.
\item D. Lopez-Paz, K. Muandet, B. Schölkopf, and I. Tolstikhin. Towards a learning theory of cause-effect inference. In Proceedings of the 32nd International Conference on Machine Learning (ICML), pages 1452--1461, 2015.
\item H. Lütkepohl. New Introduction to Multiple Time Series Analysis. Springer, Berlin, Germany, 2007.
\item M. H. Maathuis, D. Colombo, M. Kalisch, and P. Bühlmann. Estimating high-dimensional intervention effects from observational data. The Annals of Statistics, 37(6A):3133--3164, 2009.
\item D. J. C. MacKay. Information Theory, Inference, and Learning Algorithms. Cambridge University Press, New York, NY, 2002.
\item S. Mani, G. F. Cooper, and P. Spirtes. A theoretical study of y structures for causal discovery. In Proceedings of the 22nd Annual Conference on Uncertainty in Artificial Intelligence (UAI), pages 314--323, 2006.
\item D. Marinazzo, M. Pellicoro, and S. Stramaglia. Kernel method for nonlinear Granger causality. Physical Review Letters, 100:144103--1--4, 2008.
\item F. Markowetz, S. Grossmann, and R. Spang. Probabilistic soft interventions in conditional Gaussian networks. In Proceedings of the 10th International Conference on Artificial Intelligence and Statistics (AISTATS), pages 214--221, 2005.
\item J. Marschak. Statistical inference in economics: An introduction. In T. C. Koopmans, editor, Statistical Inference in Dynamic Economic Models, pages 1--50. Wiley and Chapman \& Hall, New York, NY, and London, UK, 1950. Cowles Commission for Research in Economics, Monograph No. 10.
\item B. D. McKay. Acyclic digraphs and eigenvalues of (0, 1)--matrices. Journal of Integer Sequences, 7(2):1--5, 2004.
\item C. Meek. Causal inference and causal explanation with background knowledge. In Proceedings of the 11th Annual Conference on Uncertainty in Artificial Intelligence (UAI), pages 403--441, 1995.
\item J. M. Mooij, D. Janzing, J. Peters, and B. Schölkopf. Regression by dependence minimization and its application to causal inference. In Proceedings of the 26th International Conference on Machine Learning (ICML), pages 745--752, 2009.
\item J. M. Mooij, D. Janzing, T. Heskes, and B. Schölkopf. On causal discovery with cyclic additive noise models. In Advances in Neural Information Processing Systems 24 (NIPS), pages 639--647, 2011.
\item J. M. Mooij, D. Janzing, and B. Schölkopf. From ordinary differential equations to structural causal models: The deterministic case. In Proceedings of the 29th Annual Conference on Uncertainty in Artificial Intelligence (UAI), pages 440--448, 2013.
\item J. M. Mooij, J. Peters, D. Janzing, J. Zscheischler, and B. Schölkopf. Distinguishing cause from effect using observational data: methods and benchmarks. Journal of Machine Learning Research, 17:1--102, 2016.
\item S. L. Morgan and C. Winship. Counterfactuals and Causal Inference: Methods and Principles for Social Research. Cambridge University Press, New York, NY, 2nd edition, 2007.
\item W. K. Newey. Nonparametric instrumental variables estimation. American Economic Review, 103(3):550--556, 2013.
\item J. Neyman. On the application of probability theory to agricultural experiments. Essay on principles. Section 9 (translated). Statistical Science, 5:465--480, 1923.
\item M. Nielsen and I. Chuang. Quantum Computation and Quantum Information. Cambridge University Press, New York, NY, 2000.
\item C. Nowzohour and P. Bühlmann. Score-based causal learning in additive noise models. Statistics, 50(3):471--485, 2016.
\item C. Nowzohour, M. Maathuis, and P. Bühlmann. Structure learning with bow-free acyclic path diagrams. ArXiv e-prints (1508.01717), 2015.
\item OEIS Foundation Inc. The on-line encyclopedia of integer sequences. \url{http://oeis.org/A003024}, 2017. last visit of website: 05.09.2016.
\item J. Pearl. A constraint propagation approach to probabilistic reasoning. In Proceedings of the 4th Annual Conference on Uncertainty in Artificial Intelligence (UAI), pages 31--42, 1985.
\item J. Pearl. Probabilistic Reasoning in Intelligent Systems: Networks of Plausible Inference. Morgan Kaufmann Publishers Inc., San Francisco, CA, 1988.
\item J. Pearl. Belief networks revisited. Artificial Intelligence, 59:49--56, 1993.
\item J. Pearl. Causal diagrams for empirical research. Biometrika, 82(4):669--688, 1995.
\item J. Pearl. Causality: Models, Reasoning, and Inference. Cambridge University Press, New York, NY, 2nd edition, 2009.
\item J. Pearl. Trygve Haavelmo and the emergence of causal calculus. Econometric Theory, 31:152--179, 2015.
\item J. Pearl, M. Glymour, and N. P. Jewell. Causal Inference in Statistics. Wiley, New York, NY, 2016.
\item C. S. Peirce. A theory of probable inference. In C. S. Peirce, editor, Studies in Logic by Members of the Johns Hopkins University, pages 126--181. Little, Brown, and Company, Boston, MA, 1883.
\item C. S. Peirce and J. Jastrow. On small differences in sensation. Memoirs of the National Academy of Sciences, 3:73--83, 1885.
\item E. Perkovic, J. Textor, M. Kalisch, and M. Maathuis. A complete generalized adjustment criterion. In Proceedings of the 31st Annual Conference on Uncertainty in Artificial Intelligence (UAI), pages 682--691, 2015.
\item J. Peters. Asymmetries of time series under inverting their direction. Diploma Thesis, University of Heidelberg, Heidelberg, Germany, 2008.
\item J. Peters. Restricted Structural Equation Models for Causal Inference. PhD thesis, ETH Zurich and MPI for Intelligent Systems, 2012. \url{http://dx.doi.org/10.3929/ethz-a-007597940}.
\item J. Peters. On the intersection property of conditional independence and its application to causal discovery. Journal of Causal Inference, 3:97--108, 2014.
\item J. Peters and P. Bühlmann. Identifiability of Gaussian structural equation models with equal error variances. Biometrika, 101(1):219--228, 2014.
\item J. Peters and P. Bühlmann. Structural intervention distance (SID) for evaluating causal graphs. Neural Computation, 27:771--799, 2015.
\item J. Peters, D. Janzing, A. Gretton, and B. Schölkopf. Kernel methods for detecting the direction of time series. In Proccedings of the 32nd Annual Conference of the German Classification Society (GfKl 2008), pages 1--10, 2009a.
\item J. Peters, D. Janzing, A. Gretton, and B. Schölkopf. Detecting the direction of causal time series. In Proceedings of the 26th International Conference on Machine Learning (ICML), pages 801--808, 2009b.
\item J. Peters, D. Janzing, and B. Schölkopf. Identifying cause and effect on discrete data using additive noise models. In Proceedings of the 13th International Conference on Artificial Intelligence and Statistics (AISTATS), pages 597--604, 2010.
\item J. Peters, D. Janzing, and B. Schölkopf. Causal inference on discrete data using additive noise models. IEEE Transactions on Pattern Analysis and Machine Intelligence, 33:2436--2450, 2011a.
\item J. Peters, J. M. Mooij, D. Janzing, and B. Schölkopf. Identifiability of causal graphs using functional models. In Proceedings of the 27th Annual Conference on Uncertainty in Artificial Intelligence (UAI), pages 589--598, 2011b.
\item J. Peters, D. Janzing, and B. Schölkopf. Causal inference on time series using restricted structural equation models. In Advances in Neural Information Processing Systems 26 (NIPS), pages 154--162, 2013.
\item J. Peters, J. M. Mooij, D. Janzing, and B. Schölkopf. Causal discovery with continuous additive noise models. Journal of Machine Learning Research, 15:2009--2053, 2014.
\item J. Peters, P. Bühlmann, and N. Meinshausen. Causal inference using invariant prediction: identification and confidence intervals. Journal of the Royal Statistical Society, Series B: Statistical Methodology (with discussion), 78(5):947--1012, 2016.
\item K. Peterson. Night light with sleep timer, 2005. URL \url{http://www.google.com/patents/US20050007889}. US Patent App. 10/614,245; last visit of website: 19.07.2016.
\item N. Pfister, P. Bühlmann, B. Schölkopf, and J. Peters. Kernel-based tests for joint independence. Journal of the Royal Statistical Society: Series B (to appear), 2017. doi: 10.1111/rssb.12235.
\item K. R. Popper. The Logic of Scientific Discovery. Routledge, London, 2002. 1st English Edition: 1959.
\item G. E. Quinn, C. H. Shin, M. G. Maguire, and R. A. Stone. Myopia and ambient lighting at night. Nature, 399:113--114, 1999.
\item J. Quionero-Candela, M. Sugiyama, A. Schwaighofer, and N. D. Lawrence. Dataset Shift in Machine Learning. MIT Press, Cambridge, MA, 2009.
\item R Core Team. R: A Language and Environment for Statistical Computing. R Foundation for Statistical Computing, Vienna, Austria, 2016. URL \url{http://www.R-project.org}.
\item J. Ramsey. A scalable conditional independence test for nonlinear, non-Gaussian data. ArXiv e-prints (1401.5031), 2014.
\item H. Reichenbach. The Direction of Time. University of California Press, Berkeley, CA, 1956.
\item T. S. Richardson. Markov properties for acyclic directed mixed graphs. Scandinavian Journal of Statistics, 30(1):145--157, 2003.
\item T. S. Richardson and J. M. Robins. Single world intervention graphs (SWIGs): A unification of the counterfactual and graphical approaches to causality, 2013. Working Paper Number 128, Center for Statistics and the Social Sciences, University of Washington.
\item T. S. Richardson and P. Spirtes. Ancestral graph Markov models. The Annals of Statistics, 30(4):962--1030, 2002.
\item T. S. Richardson, J. M. Robins, and I. Shpitser. Nested Markov properties for acyclic directed mixed graphs (abstract only). In Proceedings of the 28th Annual Conference on Uncertainty in Artificial Intelligence (UAI), 2012.
\item T. S. Richardson, R. J. Evans, J. M. Robins, and I. Shpitser. Nested Markov properties for acyclic directed mixed graphs. ArXiv e-prints (1701.06686), 2017.
\item J. M. Robins. A new approach to causal inference in mortality studies with sustained exposure periods --- applications to control of the healthy worker survivor effect. Mathematical Modeling, 7:1393--1512, 1986.
\item J. M. Robins, R. Scheines, P. Spirtes, and L. Wasserman. Uniform consistency in causal inference. Biometrika, 90(3):491--515, 2003.
\item R. W. Robinson. Enumeration of acyclic digraphs. In Proceedings of the 2nd Chapel Hill Conference on Combinatorial Mathematics and its Applications (University of North Carolina), pages 391--399, 1970.
\item R. W. Robinson. Counting labeled acyclic digraphs. In F. Harary, editor, New Directions in the Theory of Graphs, pages 239--273. Academic Press, New York, NY, 1973.
\item N. J. Roese. Counterfactual thinking. Psychological Bulletin, 121:133--148, 1997.
\item M. Rojas-Carulla, B. Schölkopf, R. Turner, and J. Peters. Causal transfer in machine learning. ArXiv e-prints (1507.05333v3), 2016.
\item T. Roos, T. Silander, P. Kontkanen, and P. Myllymaki. Bayesian network structure learning using factorized nml universal models. In 2008 Information Theory and Applications Workshop, pages 272--276, 2008.
\item P. R. Rosenbaum and D. B. Rubin. The central role of the propensity score in observational studies for causal effects. Biometrika, 70(1):41--55, 1983.
\item D. Rothenhäusler, C. Heinze, J. Peters, and N. Meinshausen. backShift: Learning causal cyclic graphs from unknown shift interventions. In Advances in Neural Information Processing Systems 28 (NIPS), pages 1513--1521, 2015.
\item D. B. Rubin. Estimating causal effects of treatments in randomized and nonrandomized studies. Journal of Educational Psychology, 66:688--701, 1974.
\item D. B. Rubin. Direct and indirect causal effects via potential outcomes. Scandinavian Journal of Statistics, 31(2):161--170, 2004.
\item D. B. Rubin. Causal inference using potential outcomes. Journal of the American Statistical Association, 100(469):322--331, 2005.
\item K. Sadeghi and S. Lauritzen. Markov properties for mixed graphs. Bernoulli, 20(2):676--696, 2014.
\item B. Schölkopf and A. J. Smola. Learning with Kernels. MIT Press, Cambridge, MA, 2002.
\item B. Schölkopf, D. Janzing, J. Peters, E. Sgouritsa, K. Zhang, and J. M. Mooij. On causal and anticausal learning. In Proceedings of the 29th International Conference on Machine Learning (ICML), pages 1255--1262, 2012.
\item B. Schölkopf, D. W. Hogg, D. Wang, D. Foreman-Mackey, D. Janzing, C.-J. Simon-Gabriel, and J. Peters. Removing systematic errors for exoplanet search via latent causes. In Proceedings of the 32nd International Conference on Machine Learning (ICML), pages 2218--2226, 2015.
\item B. Schölkopf, D. W. Hogg, D. Wang, D. Foreman-Mackey, D. Janzing, C.-J. Simon-Gabriel, and J. Peters. Modeling confounding by half-sibling regression. Proceedings of the National Academy of Sciences, 113(27):7391--7398, 2016.
\item T. Schreiber. Measuring information transfer. Physical Review Letters, 85:461--464, 2000.
\item E. Sgouritsa, D. Janzing, J. Peters, and B. Schölkopf. Identifying finite mixtures of nonparametric product distributions and causal inference of confounders. In Proceedings of the 29th Annual Conference on Uncertainty in Artificial Intelligence (UAI), pages 556--565, 2013.
\item E. Sgouritsa, D. Janzing, P. Hennig, and B. Schölkopf. Inference of cause and effect with unsupervised inverse regression. In Proceedings of the 18th International Conference on Artificial Intelligence and Statistics (AISTATS), pages 847--855, 2015.
\item N. Shajarisales, D. Janzing, B. Schölkopf, and M. Besserve. Telling cause from effect in deterministic linear dynamical systems. In Proceedings of the 32nd International Conference on Machine Learning (ICML), pages 285--294, 2015.
\item N. A. Sheehan, M. Bartlett, and J. Cussens. Improved maximum likelihood reconstruction of complex multi-generational pedigrees. Theoretical Population Biology, 97:11--19, 2014.
\item S. Shimizu, P. O. Hoyer, A. Hyvärinen, and A. J. Kerminen. A linear non-Gaussian acyclic model for causal discovery. Journal of Machine Learning Research, 7: 2003--2030, 2006.
\item S. Shimizu, T. Inazumi, Y. Sogawa, A. Hyvärinen, Y. Kawahara, T. Washio, P. O. Hoyer, and K. Bollen. DirectLiNGAM: A direct method for learning a linear non-Gaussian structural equation model. Journal of Machine Learning Research, 12:1225--1248, 2011.
\item H. Shimodaira. Improving predictive inference under covariate shift by weighting the log-likelihood function. Journal of Statistical Planning and Inference, 90(2):227--244, 2000.
\item A. Shojaie and G. Michailidis. Penalized likelihood methods for estimation of sparse high dimensional directed acyclic graphs. Biometrika, 97(3):519--538, 2010.
\item I. Shpitser and J. Pearl. Identification of joint interventional distributions in recursive semi-Markovian causal models. In Proceedings of the 21st AAAI Conference on Artificial Intelligence --- Volume 2, pages 1219--1226, 2006.
\item I. Shpitser and J. Pearl. Complete identification methods for the causal hierarchy. Journal of Machine Learning Research, 9:1941--1979, 2008a.
\item I. Shpitser and J. Pearl. Dormant independence. In Proceedings of the 23rd AAAI Conference on Artificial Intelligence, pages 1081--1087, 2008b.
\item I. Shpitser, T. J. VanderWeele, and J. M. Robins. On the validity of covariate adjustment for estimating causal effects. In Proceedings of the 26th Annual Conference on Uncertainty in Artificial Intelligence (UAI), pages 527--536, 2010.
\item I. Shpitser, T. S. Richardson, J. M. Robins, and R. Evans. Parameter and structure learning in nested Markov models. ArXiv e-prints (1207.5058), 2012.
\item I. Shpitser, R. J. Evans, T. S. Richardson, and J. M. Robins. Introduction to nested Markov models. Behaviormetrika, 41:3--39, 2014.
\item T. Silander and P. Myllymak. A simple approach for finding the globally optimal Bayesian network structure. In Proceedings of the 22nd Annual Conference on Uncertainty in Artificial Intelligence (UAI), pages 445--452, 2006.
\item R. Silva and R. Evans. Causal inference through a witness program. In Advances in Neural Information Processing Systems 27 (NIPS), pages 298--306, 2014.
\item R. Silva and Z. Ghahramani. The hidden life of latent variables: Bayesian learning with mixed graph models. Journal of Machine Learning Research, 10:1187--1238, 2009.
\item R. Silva, R. Scheines, C. Glymour, and P. Spirtes. Learning the structure of linear latent variable models. Journal of Machine Learning Research, 7:191--246, 2006.
\item H. A. Simon. Causal ordering and identifiability. In W. C. Hood and T. C. Koopmans, editors, Studies in Econometric Methods, pages 49--74. Wiley, New York, NY, 1953. Cowles Commission for Research in Economics, Monograph No. 14.
\item E. H. Simpson. The interpretation of interaction in contingency tables. Journal of the Royal Statistical Society, Series B: Statistical Methodology, 13:238--241, 1951.
\item V. P. Skitovič. Linear forms in independent random variables and the normal distribution law (in Russian). Izvestiia Akademii Nauk SSSR, Serija Matematiceskie, 18:185--200, 1954.
\item V. P. Skitovič. Linear combinations of independent random variables and the normal distribution law. Selected Translations in Mathematical Statistics and Probability, 2:211--228, 1962.
\item R. Solomonoff. A formal theory of inductive inference. Information and Control, Part II, 7(2):224--254, 1964.
\item C. Spearman. General intelligence, objectively determined and measured. The American Journal of Psychology, 15(2):201--292, 1904.
\item P. Spirtes. An anytime algorithm for causal inference. In Proceedings of the 8th International Conference on Artificial Intelligence and Statistics (AISTATS), pages 213--221, 2001.
\item P. Spirtes. Introduction to causal inference. Journal of Machine Learning Research, 11:1643--1662, 2010.
\item P. Spirtes, C. Glymour, and R. Scheines. Causation, Prediction, and Search. MIT Press, Cambridge, MA, 2nd edition, 2000.
\item J. Splawa-Neyman, D. M. Dabrowska, and T. P. Speed. On the application of probability theory to agricultural experiments. Essay on principles. Section 9. Statistical Science, 5(4):465--472, 1990.
\item W. Spohn. Stochastic independence, causal independence, and shieldability. Journal of Philosophical Logic, 9:73--99, 1980.
\item R. P. Stanley. Acyclic orientations of graphs. Discrete Mathematics, 7(5):171--178, 1973.
\item I. Steinwart and A. Christmann. Support Vector Machines. Springer, New York, NY, 2008.
\item B. Steudel and N. Ay. Information-theoretic inference of common ancestors. Entropy, 17(4):2304--2327, 2015.
\item B. Steudel, D. Janzing, and B. Schölkopf. Causal Markov condition for submodular information measures. In Proceedings of the 23rd Annual Conference on Learning Theory (COLT), pages 464--476, 2010.
\item M. Studený and D. Haws. Learning Bayesian network structure: Towards the essential graph by integer linear programming tools. International Journal of Approximate Reasoning, 55:1043--1071, 2014.
\item M. Sugiyama and M. Kawanabe. Machine Learning in Non-Stationary Environment. MIT Press, Cambridge, MA, 2012.
\item X. Sun, D. Janzing, and B. Schölkopf. Causal inference by choosing graphs with most plausible Markov kernels. In Proceedings of the 9th International Symposium on Artificial Intelligence and Mathematics, pages 1--11, 2006.
\item R. S. Sutton and A. G. Barto. Reinforcement Learning: An Introduction. MIT Press, Cambridge, MA, 2nd edition, 2015.
\item M. Teyssier and D. Koller. Ordering-based search: A simple and effective algorithm for learning Bayesian networks. In Proceedings of the 21st Annual Conference on Uncertainty in Artificial Intelligence (UAI), pages 584--590, 2005.
\item J. Tian. Studies in Causal Reasoning and Learning. PhD thesis, Department of Computer Science, University of California, Los Angeles, CA, 2002.
\item J. Tian and J. Pearl. Causal discovery from changes. In Proceedings of the 17th Annual Conference on Uncertainty in Artificial Intelligence (UAI), pages 512--522, 2001.
\item J. Tian and J. Pearl. On the testable implications of causal models with hidden variables. In Proceedings of the 18th Annual Conference on Uncertainty in Artificial Intelligence (UAI), pages 519--527, 2002.
\item R. Tillman, A. Gretton, and P. Spirtes. Nonlinear directed acyclic structure learning with weakly additive noise models. In Advances in Neural Information Processing Systems 22 (NIPS), pages 1847--1855, 2009.
\item R. E. Tillman and F. Eberhardt. Learning causal structure from multiple datasets with similar variable sets. Behaviormetrika, 41(1):41--64, 2014.
\item S. Triantafillou and I. Tsamardinos. Constraint-based causal discovery from multiple interventions over overlapping variable sets. Journal of Machine Learning Research, 16:2147--2205, 2015.
\item S. Triantafillou, I. Tsamardinos, and I. G. Tollis. Learning causal structure from overlapping variable sets. In Proceedings of the 13th International Conference on Artificial Intelligence and Statistics (AISTATS), pages 860--867, 2010.
\item I. Tsamardinos, L. E. Brown, and C. F. Aliferis. The max-min hill-climbing Bayesian network structure learning algorithm. Machine Learning, 65(1):31--78, 2006.
\item C. Uhler, G. Raskutti, P. Bühlmann, and B. Yu. Geometry of the faithfulness assumption in causal inference. The Annals of Statistics, 41(2):436--463, 2013.
\item S. Ullman. The Interpretation of Visual Motion. MIT Press, Cambridge, MA, 1979.
\item R. Urner, S. Shalev-Shwartz, and S. Ben-David. Access to unlabeled data can speed up prediction time. In Proceedings of the 28th International Conference on Machine Learning (ICML), pages 641--648, 2011.
\item S. A. van de Geer. Empirical Processes in M-Estimation. Cambridge Series in Statistical and Probabilistic Mathematics. Cambridge University Press, Cambridge, UK, 2009.
\item V. N. Vapnik. Statistical Learning Theory. Wiley, New York, NY, 1998.
\item T. Verma and J. Pearl. Causal networks: Semantics and expressiveness. In Proceedings of the 4th Annual Conference on Uncertainty in Artificial Intelligence (UAI), pages 352--359, 1988.
\item T. Verma and J. Pearl. Equivalence and synthesis of causal models. In Proceedings of the 6th Annual Conference on Uncertainty in Artificial Intelligence (UAI), pages 255--270, 1991.
\item D. Voiculescu, editor. Free Probability Theory, volume 12 of Fields Institute Communications. American Mathematical Society, Providence, RI, 1997.
\item A. Wald. Statistical Decision Functions. Wiley, New York, NY, 1950.
\item L. Wang and E. Tchetgen Tchetgen. Bounded, efficient and triply robust estimation of average treatment effects using instrumental variables. ArXiv e-prints (1611.09925), 2016.
\item H. White and X. Lu. Granger causality and dynamic structural systems. Journal of Financial Econometrics, 8(2):193--243, 2010.
\item N. Wiener. The theory of prediction. In E. Beckenbach, editor, Modern Mathematics for Engineers. McGraw-Hill, New York, NY, 1956.
\item H. P. Williams. Fourier's method of linear programming and its dual. The American Mathematical Monthly, 93(9):681--695, 1986.
\item J. Wishart. Sampling errors in the theory of two factors. British Journal of Psychology, pages 180--187, 1928.
\item S. N. Wood. Generalized Additive Models: An Introduction with R. Chapman \& Hall/CRC, London, UK, 2006.
\item P. G. Wright. The Tariff on Animal and Vegetable Oils. Investigations in International Commercial Policies. Macmillan, New York, NY, 1928.
\item S. Wright. On the nature of size factors. Genetics, 3:367--374, 1918.
\item S. Wright. The relative importance of heredity and environment in determining the piebald pattern of guinea-pigs. Proceedings of the National Academy of Sciences, 6(6):320--332, 1920.
\item S. Wright. Correlation and causation. Journal of Agricultural Research, 20(7): 557--585, 1921.
\item S. Wright. The method of path coefficients. Annals of Mathematical Statistics, 5(3):161--215, 1934.
\item K. Zadnik, L. A. Jones, B. C. Irvin, R. N. Kleinstein, R. E. Manny, J. A. Shin, and D. O. Mutti. Vision: Myopia and ambient night-time lighting. Nature, 404: 143--144, 2000.
\item J. Zhang. On the completeness of orientation rules for causal discovery in the presence of latent confounders and selection bias. Artificial Intelligence, 172: 1873--1896, 2008a.
\item J. Zhang. Causal reasoning with ancestral graphs. Journal of Machine Learning Research, 9:1437--1474, 2008b.
\item J. Zhang and P. Spirtes. Strong faithfulness and uniform consistency in causal inference. In Proceedings of the 19th Annual Conference on Uncertainty in Artificial Intelligence (UAI), pages 632--639, 2003.
\item J. Zhang and P. Spirtes. A characterization of Markov equivalence classes for ancestral graphical models. Technical Report No. CMU-PHIL-168, 2005.
\item J. Zhang and P. Spirtes. Detection of unfaithfulness and robust causal inference. Minds and Machines, 18(2):239--271, 2008.
\item K. Zhang and L. Chan. Extensions of ICA for causality discovery in the Hong Kong stock market. In 13th International Conference on Neural Information Processing, (ICONIP), pages 400--409, 2006.
\item K. Zhang and A. Hyvärinen. On the identifiability of the post-nonlinear causal model. In Proceedings of the 25th Annual Conference on Uncertainty in Artificial Intelligence (UAI), pages 647--655, 2009.
\item K. Zhang, J. Peters, D. Janzing, and B. Schölkopf. Kernel-based conditional independence test and application in causal discovery. In Proceedings of the 27th Annual Conference on Uncertainty in Artificial Intelligence (UAI), pages 804--813, 2011.
\item K. Zhang, B. Schölkopf, K. Muandet, and Z. Wang. Domain adaptation under target and conditional shift. In Proceedings of the 30th International Conference on Machine Learning (ICML), pages 819--827, 2013.
\item K. Zhang, M. Gong, and B. Schölkopf. Multi-source domain adaptation: A causal view. In Proceedings of the 29th AAAI Conference on Artificial Intelligence, pages 3150--3157, 2015.
\item J. Zscheischler, D. Janzing, and K. Zhang. Testing whether linear equations are causal: A free probability theory approach. In Proceedings of the 27th Annual Conference on Uncertainty in Artificial Intelligence (UAI), pages 839--846, 2011.
\item W. Zurek. Algorithmic randomness and physical entropy. Physical Review A, 40(8):4731--4751, 1989.
\end{list}
\endgroup
